\documentclass[12pt]{article}
\usepackage[a4paper, margin=1in]{geometry}
\usepackage{enumitem}
\usepackage{setspace}

\setstretch{1.2}

\begin{document}

\begin{center}
    {\LARGE \textbf{Lecture 23 Scribe}}\\[0.3cm]
    {\large \textbf{Predicting the Impossible: Tail Bounds and System Reliability}}\\[0.2cm]
    {\normalsize \textit{Case Study: Amazon Black Friday Sale 2026 (Coding Session – Tails and Bounds)}}
\end{center}

\vspace{0.5cm}

\section*{1. Introduction}

\begin{itemize}
    \item \textbf{Lecture Title:} Predicting the Impossible: Tail Bounds and System Reliability  
    \item \textbf{Context:} Case Study on Amazon Black Friday Sale 2026  
    \item \textbf{Instructor:} Dhaval Patel, PhD  
    \item \textbf{Affiliation:} Computer Science and Engineering (CSE), Ahmedabad University  
    \item \textbf{Date:} April 02, 2026  
\end{itemize}

\section*{2. Case Study: Amazon Black Friday Sale 2026}

\begin{itemize}
    \item The \textbf{Amazon Black Friday Sale 2026} is expected to begin on:
    \begin{itemize}
        \item \textbf{Friday, November 27, 2026}
    \end{itemize}
    \item The timeline is based on:
    \begin{itemize}
        \item Historical patterns  
        \item Forecasts from external sources (e.g., 91mobiles, Bajaj Finserv)
    \end{itemize}
\end{itemize}

\subsection*{Purpose of Case Study}

\begin{itemize}
    \item To connect:
    \begin{itemize}
        \item \textbf{Uncertainty}
        \item \textbf{Computing}
    \end{itemize}
    \item The system under consideration involves:
    \begin{itemize}
        \item Large-scale user participation  
        \item High demand variability  
        \item Uncertain system behavior  
    \end{itemize}
\end{itemize}

\section*{3. Conceptual Framework: Uncertainty and Computing}

\begin{itemize}
    \item The lecture establishes a connection between:
    \begin{itemize}
        \item Real-world large-scale systems (e.g., online sale events)
        \item Mathematical tools for handling uncertainty  
    \end{itemize}
    \item The central problem:
    \begin{itemize}
        \item \textbf{Predicting rare or extreme events} in system performance  
    \end{itemize}
\end{itemize}

\section*{4. Tail Bounds}

Tail bounds are introduced as mathematical tools to quantify the probability of extreme deviations.

\subsection*{4.1 Markov’s Inequality}

\begin{itemize}
    \item Introduced as a fundamental inequality for bounding probabilities  
    \item Applies to:
    \begin{itemize}
        \item Non-negative random variables  
    \end{itemize}
\end{itemize}

\subsection*{4.2 Chebyshev’s Inequality}

\begin{itemize}
    \item Extends the idea of bounding probabilities  
    \item Incorporates:
    \begin{itemize}
        \item Mean  
        \item Variance  
    \end{itemize}
\end{itemize}

\subsection*{4.3 Chernoff Bound}

\begin{itemize}
    \item Provides tighter bounds compared to:
    \begin{itemize}
        \item Markov’s Inequality  
        \item Chebyshev’s Inequality  
    \end{itemize}
    \item Used for:
    \begin{itemize}
        \item Bounding tail probabilities of sums of random variables  
    \end{itemize}
\end{itemize}

\section*{5. Application to System Reliability}

\begin{itemize}
    \item Tail bounds are used to:
    \begin{itemize}
        \item Analyze \textbf{system reliability}
        \item Estimate \textbf{probability of failure or overload}
    \end{itemize}
    \item In the context of the Amazon sale:
    \begin{itemize}
        \item System must handle:
        \begin{itemize}
            \item Sudden spikes in traffic  
            \item High uncertainty in demand  
        \end{itemize}
    \end{itemize}
\end{itemize}

\section*{6. Lecture Flow and Structure}

\begin{itemize}
    \item The lecture progresses through:
    \begin{enumerate}
        \item Introduction to real-world uncertainty (Amazon sale case)
        \item Motivation for probabilistic bounds  
        \item Presentation of three inequalities:
        \begin{itemize}
            \item Markov’s Inequality  
            \item Chebyshev’s Inequality  
            \item Chernoff Bound  
        \end{itemize}
        \item Connection to system reliability  
    \end{enumerate}
\end{itemize}

\section*{7. Class Participation}

\begin{itemize}
    \item Slides indicate:
    \begin{itemize}
        \item Active engagement during the lecture  
        \item Participation-based interaction  
    \end{itemize}
\end{itemize}

\section*{8. Key Logical Dependencies}

\begin{itemize}
    \item Real-world uncertainty $\rightarrow$ Need for probabilistic modeling  
    \item Probabilistic modeling $\rightarrow$ Use of tail bounds  
    \item Tail bounds $\rightarrow$ Quantification of extreme events  
    \item Quantification $\rightarrow$ Application to system reliability  
\end{itemize}

\section*{9. Summary of Included Content}

\begin{itemize}
    \item Definitions:
    \begin{itemize}
        \item Markov’s Inequality  
        \item Chebyshev’s Inequality  
        \item Chernoff Bound  
    \end{itemize}
    \item Context:
    \begin{itemize}
        \item Amazon Black Friday Sale 2026  
    \end{itemize}
    \item Conceptual Link:
    \begin{itemize}
        \item Uncertainty $\leftrightarrow$ Computing  
    \end{itemize}
    \item Application:
    \begin{itemize}
        \item System reliability under uncertain demand  
    \end{itemize}
\end{itemize}

\vspace{0.5cm}

\noindent \textit{End of Scribe}

\end{document}